% !PRINTED 311
\weberpage{311}
%<<CONT
der Anfang der Entwickelung für $w_z$:
\weq{4}{w_z = q^{-\frac{2}{5}} e^{\frac{c}{2}\frac{\pi i}{5}} + \cdots}
Hieraus folgt, daß die Potenzsummen der $w_z$ bis zur vierten einschließlich Konstanten sind, da sie nicht einmal die erste Potenz von $j(\omega)$ enthalten können, und daß das Produkt der $w_z$ eine \gesperrt{lineare} Funktion von $j(w)$ sein muß, % [sic] Gedruckt ist w; vier Wörter später steht ω.
in der $j(\omega)$ den Koeffizienten 1 hat.

Die Gleichung der $w$ hat daher die Form:
\weq{5}{w^5 + b_1 w^4 + b_2 w^3 + b_3 w^2 + b_4 w + b_5 = j(\omega),}
wenn die $b$ rationale Zahlen sind. Diese Koeffizienten lassen sich dadurch berechnen, daß man die Entwickelung (4) weiter fortsetzt. Wir schlagen hier einen anderen Weg ein, der, streng genommen, zu der im nächsten Teil behandelten komplexen Multiplikation gehört, bei seiner Einfachheit aber trotzdem ganz wohl hier seine Stelle finden kann.

Setzen wir $\omega = i = \sqrt{-1}$, so ist
\[ \omega = -\frac{1}{\omega} \]
und infolgedessen ist [§ 34, (11), (12), (14)]:
\[ f_1(i) = f_2(i) = \sqrt[8]{2}, \quad f(i) = \sqrt[4]{2}, \]
also [§ 54, (5)]:
\weq{6}{\gamma_3(i) = 0, \quad \gamma_2(i) = 12, \quad j(i) = 12^3.}
Man findet aber ferner aus den Transformationsformeln für $\eta(\omega)$ [§ 34, (4), (5)]:
\[ \eta(\omega+1) = e^{\frac{\pi i}{12}} \eta(\omega), \quad \eta\left(-\frac{1}{\omega}\right) = \sqrt{-i\omega}\,\eta(\omega), \]
wenn man auf die Zerlegung
\[ 5 = (2+i)(2-i) \]
Rücksicht nimmt:
\[ \eta\left(\frac{12+i}{5}\right) = e^{\frac{\pi i}{6}} \eta\left(\frac{2+i}{5}\right) = e^{\frac{\pi i}{6}} \eta\left(\frac{1}{2-i}\right) = \sqrt{1+2i}\,\eta(i), \]
also
\[ v_2 = 1+2i, \quad v_3 = 1-2i, \]
und man kann jetzt, da für $\gamma_2 = 12$ zwei Wurzeln der Gleichung (1) bekannt sind, für diesen besonderen Wert von $\gamma_2$ den linken Teil dieser Gleichung leicht in Faktoren zerlegen:
\[ v^6 + 10v^3 - 12v + 5 = (v^2+2v+5)(v^2+v-1)^2, \]
%CONT>>
